\chapter{Results}
\label{ch:results}
The results are organised by the question each experiment can answer.
Normal profiles describe the fitted workloads. Development comparisons
explain the adaptation and scoring choices. Numerical checks compare
implementations of the same objective. The duration study examines a
different observation scale, while the closing replay records the selected
configuration on its designated final partition. The last sections report
the constrained runtime and software verification. Tables and plot data
are generated from the retained artifacts rather than transcribed from
earlier narrative summaries.

\section{Normal traffic profiles and static development detection}
\label{sec:normal-results}

\subsection{Different devices require different reference mixtures}

Table~\ref{tab:profiles} reports the fitted mean proportions and total
concentration. The D-Link model distributes substantial mass across all
four categories. Chromecast is dominated by ordinary UDP, with an
expected share of about 92.6\%; its SSDP parameter is positive but very
small. Samsung assigns most expected mass to SSDP and TCP, with little
ordinary UDP. Figure~\ref{fig:profiles} displays these contrasts on the
same scale. The profiles are properties of the selected normal captures,
not universal descriptions of the commercial device models.

\begin{table}[tbp]
\centering\small
\begin{tabular}{@{}lrrrrrr@{}}
\toprule
Device & TCP (\%) & UDP (\%) & SSDP (\%) & ARP (\%) & $\alpha_0$ & $\psi$ \\
\midrule
\TableRows{generated/profiles}
\bottomrule
\end{tabular}
\caption{Saved normal-model parameters expressed as expected category shares,
total concentration $\alpha_0$, and inverse concentration $\psi=1/\alpha_0$.
UDP excludes SSDP. Shares are rounded independently.}
\label{tab:profiles}
\end{table}

\begin{figure}[tbp]
\centering
\begin{tikzpicture}
\begin{axis}[width=.95\textwidth,height=7.1cm,ybar,bar width=9pt,
 ymin=0,ymax=105,ylabel={Fitted mean share (\%)},
 xtick={1,2,3,4},xticklabels={TCP,Ordinary UDP,SSDP,ARP},
 legend style={at={(.5,1.03)},anchor=south,legend columns=3,font=\small,column sep=6pt},
 ymajorgrids=true,grid style={gray!20}]
\addplot[fill=DLinkGreen!60,draw=DLinkGreen] table[x=category,y=percent,col sep=comma]{generated/profile_D-LinkDayCam5.csv};
\addplot[fill=ChromecastBlue!60,draw=ChromecastBlue] table[x=category,y=percent,col sep=comma]{generated/profile_UNSW-Chromecast.csv};
\addplot[fill=SamsungOrange!60,draw=SamsungOrange] table[x=category,y=percent,col sep=comma]{generated/profile_UNSW-Samsung-Camera.csv};
\legend{D-Link,Chromecast,Samsung}
\end{axis}
\end{tikzpicture}
\caption{Expected category proportions, $p_k=\alpha_k/\alpha_0$, from the saved
normal models. The bars are fitted means, not pooled packet fractions or measured
attack signatures. Dataset-specific models have different expected mixtures.}
\label{fig:profiles}
\end{figure}


These differences matter when interpreting a count anomaly. An SSDP-heavy
window is not inherently anomalous: it is far less surprising under the
Samsung reference than under the Chromecast reference. Likewise, an ARP
count must be compared with the selected device's benign scale. A shared
feature vocabulary therefore does not imply a shared fitted distribution
or a globally meaningful raw likelihood cutoff.

The fitted concentration also differs substantially. D-Link has
$\alpha_0\simeq122.2$, Chromecast $28.1$, and Samsung $17.4$. For a fixed
total count, the lower-concentration models permit more variation in
composition around their means. This is the model's overdispersion
parameter, rather than a measure of attack prevalence or confidence in
the training labels. The larger D-Link concentration cannot establish
that D-Link traffic is more secure.

Table~\ref{tab:fit-parameters} records the individual parameters, iteration
counts, and stored DM calibration quantities. Chromecast's small positive
SSDP parameter illustrates why zero category counts and positive parameter
validation need careful numerical treatment. The three fits converge in
different numbers of iterations even though the same tolerance and
maximum-iteration setting are used. A compact four-category model is not
necessarily a one-step fit.

\begin{table}[tbp]
\centering\footnotesize
\begin{tabular}{@{}lrrrrrrr@{}}
\toprule
Device & $\alpha_{TCP}$ & $\alpha_{UDP}$ & $\alpha_{SSDP}$ & $\alpha_{ARP}$ & Iter. & $\tau$ & $d_1$ \\
\midrule
\TableRows{generated/parameters}
\bottomrule
\end{tabular}
\caption{Saved LM fits and current joint-calibration DM cutoffs and scales.
The cutoff $\tau$ is in complete log-probability units, while $d_1$ is the
stored fitting-score scale. These scales should not be substituted for the
different historical adaptation scales.}
\label{tab:fit-parameters}
\end{table}

\subsection{Static decisions on labelled development windows}

The reference static detector produces the counts in
Table~\ref{tab:development}. Chromecast detects 15 of 32 positive windows
and Samsung 48 of 70. The corresponding F1 values are 61.22\% and
78.69\%. Both models produce relatively few false alerts on these
development observations: two among 442 benign Chromecast windows and
four among 486 benign Samsung windows. The false positive rates are
approximately 0.45\% and 0.82\%, respectively.

\begin{table}[tbp]
\centering\small
\begin{tabular}{@{}lrrrrrrr@{}}
\toprule
Device & TP & FP & FN & TN & Precision & Recall & F1 \\
\midrule
\TableRows{generated/development}
\bottomrule
\end{tabular}
\caption{Reference static development results. The retained current joint
rule and the historical static DM rule give the same binary counts on these
observations. This equality does not extend to every adaptive comparison.}
\label{tab:development}
\end{table}

The higher precision than recall means that a large part of the remaining
error consists of attacks that do not cross the chosen normal-data
cutoffs. This is especially clear for Chromecast, whose development
recall is 46.88\%. Neither normal-only fitting nor the one-percent
calibration quantile optimises recall against these attack labels.
A benign tail setting must not be interpreted as a guarantee of
detecting 99\% of attacks.

The D-Link development workload produces three alerts in 270 windows,
and its later numerical partition produces six in 288. Without attack
annotations these are alert counts only. Calling all nine false
positives would add an assumption absent from the data. The labelled
devices therefore carry the detection conclusions, while D-Link
primarily confirms that the numerical comparison exercises both
alerting and non-alerting decisions.

\subsection{Joint calibration produces unequal count sensitivities}

Table~\ref{tab:count-cutoffs} expresses the count-tail cutoffs both in
their stored log space and as equivalent packet-count boundaries.
For a total count, the comparison is $\log(1+N)>u_N$; the boundary
is therefore $\exp(u_N)-1$. Values in the count columns are rounded
for interpretation, while decisions use the unrounded log cutoffs.

\begin{table}[tbp]
\centering\small
\begin{tabular}{@{}lrrrr@{}}
\toprule
Device & $u_N$ & Total boundary & $u_A$ & ARP boundary \\
\midrule
\TableRows{generated/count_cutoffs}
\bottomrule
\end{tabular}
\caption{Saved count-tail cutoffs and rounded equivalent packet-count
boundaries. These are detection cutoffs; the separate learning-volume cap
is a different empirical quantity.}
\label{tab:count-cutoffs}
\end{table}

The differences are substantial. Samsung's total boundary is about
1,610 packets, while Chromecast's is about $1.12\times10^{13}$.
The latter total-volume guard is consequently inactive for ordinary
count ranges in the evaluated traffic. Its presence in the decision
rule does not mean that it contributes useful sensitivity on every
device. The actual case margins and static-equivalence results are
consistent with this calibration outcome.

This follows from calibrating the minimum of standardised features
with one common tail quantile. A sufficiently extreme benign tail
in one feature can require a negative joint quantile that also
makes another feature's implied cutoff very permissive. Median/IQR
scaling supplies units for combining the features; it does not
equalise their lower tails or guarantee that they all trigger.
The saved cutoffs therefore need to be inspected alongside the
formula, rather than interpreting the count checks as three
independently active protections.

The result also separates detection from learning admission. The
conservative buffer uses a distinct 99th-percentile calibration
volume cap. Even if the joint total guard is permissive, that
admission cap can still reject high-volume learning observations.
Neither cutoff is attack-labelled, and their different purposes
should remain visible when analysing an adaptive run.

\section{Historical adaptation and parameter movement}
\label{sec:legacy-results}

\subsection{Development gains depend on device and schedule}

The historical policy admits windows using less restrictive checks and
allows larger parameter movement than the conservative retained policy.
Table~\ref{tab:legacy} and Figure~\ref{fig:legacy} show its development
comparisons under DM-only scoring. The names 144/144 and 288/240 specify
the attempt cadence and minimum accepted-window count. They refer to
different configurations, not successive states of one detector.

\begin{table}[tbp]
\centering\small
\begin{tabular}{@{}llrrrrr@{}}
\toprule
Device & Configuration & TP & FP & FN & TN & F1 \\
\midrule
\TableRows{generated/legacy}
\bottomrule
\end{tabular}
\caption{Historical development comparisons with a 288-window learning buffer.
The refit 288/240 configuration was selected after development inspection;
the table is not evidence from an untouched selection set.}
\label{tab:legacy}
\end{table}

\begin{figure}[tbp]
\centering
\begin{tikzpicture}
\begin{axis}[width=.95\textwidth,height=7cm,ymin=0,ymax=100,ylabel={Development F1 (\%)},
 xtick={0,1,2,3},xticklabels={Static,Threshold 144,Refit 144/144,Refit 288/240},
 x tick label style={rotate=15,anchor=east,font=\small},
 legend style={at={(.5,1.02)},anchor=south,legend columns=2,column sep=6pt},
 ymajorgrids=true,grid style={gray!20}]
\addplot[ChromecastBlue,thick,mark=*] table[x=mode,y=f1,col sep=comma]{generated/legacy_UNSW-Chromecast.csv};
\addplot[SamsungOrange,thick,dashed,mark=square*] table[x=mode,y=f1,col sep=comma]{generated/legacy_UNSW-Samsung-Camera.csv};
\legend{Chromecast,Samsung camera}
\end{axis}
\end{tikzpicture}
\caption{Historical DM-only development comparisons. Refit labels give attempt
cadence/minimum accepted windows; every adaptive variant has a 288-window buffer.
The lines connect distinct configurations, not a temporal learning trajectory.
These development-selected runs differ from the conservative closing replay.}
\label{fig:legacy}
\end{figure}


On Chromecast, the 144-window threshold policy lowers F1 from 61.22\% to
55.32\%. Refitting at 144/144 gives 58.33\%, also below static operation.
The later 288/240 refit detects 23 of 32 positive windows, with two
false positives, and reaches 80.70\% F1. The improvement comes from
eight additional true positives at the same observed false positive
count, compared with the static reference.

Samsung behaves differently. Both 144-window adaptive variants produce
27 true positives, three false positives, and 43 false negatives,
reducing F1 to 54.00\%. The selected 288/240 refit improves upon those
adaptive results, but its 77.31\% F1 remains below the static 78.69\%.
Its one fewer false positive is accompanied by two fewer true positives.
Thus the historical study does not show a consistent advantage for
refitting, even after the more favourable configuration is selected.

This comparison demonstrates why a single pooled improvement would be
incomplete. Chromecast's gain and Samsung's failure are different
responses to the same general learning mechanism. A useful policy
must control changes relative to each device's reference, rather than
relying on an average of opposing outcomes. The development results
motivate that requirement; they do not validate a particular guard
on future unseen drift.

\subsection{Concentration changes reveal what the refit learned}

Table~\ref{tab:concentration} reports initial and final concentrations.
Chromecast's changes are moderate, about 1.28 and 1.23 times its
original concentration. Samsung's are much larger: about 11.87 and
9.58 times the original value. These are substantive distribution
changes, even when the expected proportions remain recognisable.

\begin{table}[tbp]
\centering\small
\begin{tabular}{@{}llrrr@{}}
\toprule
Device & Configuration & Initial $\alpha_0$ & Final $\alpha_0$ & Ratio \\
\midrule
\TableRows{generated/concentration}
\bottomrule
\end{tabular}
\caption{Concentration movement in historical refitting. Values come from the
retained adaptation summaries, independently of classification metrics.}
\label{tab:concentration}
\end{table}

Increasing concentration reduces the allowed compositional dispersion
conditional on total count. It can change both benign tail probabilities
and attack scores without supplying a new attack feature. It would
therefore be misleading to describe the Samsung refit simply as learning
the same normal behaviour more accurately. Its observed movement is
large and its detection result is worse than the static reference.

These updates illustrate a limitation of self-selected learning data.
A window passing the active cutoff is not known to be benign. Accepted
windows are selected partly because they look compatible with the
current model, and that selection can narrow the apparent distribution.
Attack-like traffic remaining inside the cutoff can enter the same
buffer. These mechanisms are compatible with the observed movement,
but concentration alone cannot identify how much is caused by
contamination, benign drift, or selection bias. It is a diagnostic
of the update, not an attack label.

The conservative policy consequently bounds concentration, mean
proportions, and threshold movement and retains the original reference.
Later results show that this protection often prevents activation
entirely. That outcome is candidate rejection, rather than a successful
refit merely because an adaptive mode was enabled.

\section{Retrospective selection of the joint decision rule}
\label{sec:joint-results}

The count-tail study compares baseline DM scoring with the retained
joint composition, total-volume, and ARP decision rule. Its pooled
results in Table~\ref{tab:joint-comparison} cover 2,038 unique labelled
development and later-audit windows across the two devices. This is a
retrospective comparison: the later labels had already been inspected.
It explains the rule selection but does not provide an independent
estimate of generalisation.

\begin{table}[tbp]
\centering\small
\begin{tabular}{@{}lrrrrr@{}}
\toprule
Configuration & DM F1 & Joint F1 & Joint TP & Joint FP & Joint FN \\
\midrule
\TableRows{generated/scoring_comparison}
\bottomrule
\end{tabular}
\caption{Retrospective pooled scoring comparison. Baseline and joint variants
share each configuration's update policy and input support. Legacy refitting
here uses 288/144, distinct from the historical 288/240 study.}
\label{tab:joint-comparison}
\end{table}

For static operation and both conservative modes, the DM and joint
variants produce the same counts: 77 true positives, nine false
positives, 43 false negatives, and 1,909 true negatives. F1 is
74.76\%. The extra tails do not change these static decisions.
This does not establish their redundancy for every possible count
vector.

The effect differs under the legacy adaptive policy. Threshold
adaptation has 56.50\% baseline F1 and 74.76\% joint F1. Legacy
refitting has 57.30\% baseline F1 and 75.36\% joint F1, with 78 true
positives rather than the static 77. Original count references
preserve alerts that can disappear when an adapting DM cutoff or model
moves. The large gain in these rows is primarily protection against
a weakened adaptive decision, rather than a large increase over the
original static detector.

The scoring study rejected broader alternatives. Dividing a score
by packet count changed the meaning of the complete probability and
lost attack recall. Broad category-specific tails introduced later-audit
false alerts. Separate rank-based tails did not retain the desired
recall. These outcomes support the narrower rule within the inspected
data, but do not constitute a comprehensive feature-selection
benchmark. Richer features were not available uniformly in the
already prepared four-count records.

The implementation subsequently evaluates cheap count guards before
computing the DM likelihood. A regression study compares this ordering
with the earlier full-score implementation in 20 replays covering
all 2,038 labelled windows. Binary decisions, accepted-update counts,
final fingerprints, thresholds, and buffer sizes agree. Event scores
and severities may differ when an early count violation makes further
numerical evaluation unnecessary. This thesis therefore uses confusion
counts for cross-version equivalence and does not silently treat all
saved rank scores as interchangeable.

\section{LM and SciPy: numerical agreement and computational cost}
\label{sec:numerical-results}

\subsection{Agreement on the production objective}

The D-Link production benchmark fits the same 846-row, four-category
matrix with both backends. LM converges in 1,882 iterations and SciPy
in 1,879. The largest relative difference between fitted parameter
components is approximately $3.86\times10^{-7}$. The largest absolute
window-score difference in this benchmark is $1.09\times10^{-5}$,
and the calibrated cutoff differs by approximately
$2.89\times10^{-6}$. Both backends flag the same three of 270
development windows.

Different iteration counts are compatible with numerical agreement.
Small differences in intermediate special-function values can change
when a finite convergence tolerance is first satisfied. The relevant
checks are the final fitted quantities and decisions, rather than a
requirement that all iteration counters match. Conversely, matching
decisions alone would be too weak to identify a large hidden parameter
discrepancy. Both checks are reported.

The three-dataset check in Table~\ref{tab:backend-agreement} extends the
comparison to Chromecast and Samsung. Across 1,300 development windows,
no binary decision differs. The largest relative parameter difference
is about $1.71\times10^{-6}$, on Chromecast; the largest absolute
score difference is about $2.40\times10^{-4}$. The D-Link score
difference in this second record is smaller than in the production
benchmark because these are distinct saved runs, not repeated
statistics merged into one experiment.

\begin{table}[tbp]
\centering\small
\begin{tabular}{@{}lrrr@{}}
\toprule
Device & Max. relative $\alpha$ difference & Max. score difference & Decision differences \\
\midrule
\TableRows{generated/backend_agreement}
\bottomrule
\end{tabular}
\caption{LM--SciPy agreement in the three-dataset backend check. Score
differences are absolute differences in complete log-probability units;
zero decision differences is a finite-workload result.}
\label{tab:backend-agreement}
\end{table}

These results support interchangeability of the observed binary
decisions under the tested inputs and cutoffs. They do not prove
bit-for-bit equivalence for arbitrary parameters or a window arbitrarily
close to a cutoff. The score definition, coefficient, and fallback
behaviour remain part of the comparison. A stable calculation of the
wrong objective would not be validated by agreement on a few examples.

\subsection{The faster backend depends on the operation}

Table~\ref{tab:backend-timing} separates complete fitting, one matrix
likelihood call, and one window score. SciPy is faster for fitting
and the matrix kernel, while LM is faster for scoring one window.
The LM/SciPy ratios are approximately 1.064, 1.110, and 0.899.
The result is mixed rather than a uniform speed advantage.

\begin{table}[tbp]
\centering\small
\begin{tabular}{@{}lrrrl@{}}
\toprule
Operation & LM & SciPy & LM/SciPy & Unit \\
\midrule
\TableRows{generated/backend_timing}
\bottomrule
\end{tabular}
\caption{D-Link production benchmark medians. Fits include estimator iterations;
matrix and window rows measure their specified kernel boundaries after warm-up.
Lower runtime is better.}
\label{tab:backend-timing}
\end{table}

Figure~\ref{fig:backend-ratios} shows the separate three-dataset timing
check. All datasets favour LM for the single-window operation,
but favour SciPy for the matrix kernel and complete fit. The size of
the difference depends on parameters, matrix, and call boundary.
These are measurements from one environment and finite repeated
batches, not estimates for all hardware or compiler combinations.

\begin{figure}[tbp]
\centering
\begin{tikzpicture}
\begin{axis}[width=.95\textwidth,height=7cm,ymin=.7,ymax=1.85,
 ylabel={LM/SciPy median runtime ratio},
 xtick={1,2,3},xticklabels={Complete fit,Matrix kernel,Window score},
 legend style={at={(.5,1.03)},anchor=south,legend columns=3,font=\small,column sep=6pt},
 ymajorgrids=true,grid style={gray!20}]
\addplot[DLinkGreen,mark=triangle*,thick] table[x=operation,y=ratio,col sep=comma]{generated/backend_D-LinkDayCam5.csv};
\addplot[ChromecastBlue,mark=*,thick] table[x=operation,y=ratio,col sep=comma]{generated/backend_UNSW-Chromecast.csv};
\addplot[SamsungOrange,mark=square*,dashed,thick] table[x=operation,y=ratio,col sep=comma]{generated/backend_UNSW-Samsung-Camera.csv};
\addplot[black,dotted,domain=.8:3.2,forget plot]{1};
\legend{D-Link,Chromecast,Samsung}
\end{axis}
\end{tikzpicture}
\caption{Operation-specific runtime ratios from the retained three-dataset
backend-check record. Values below one favour LM. This record is distinct from
the separate D-Link production timing table; ratios from different runs are not
combined into a single speedup estimate.}
\label{fig:backend-ratios}
\end{figure}


An efficient underlying likelihood method can still incur a slower
application matrix call. The production interface preserves row-wise
probabilities, so native setup and dispatch are paid at a finer
granularity than for a pooled objective. NumPy/SciPy can evaluate
parts of the matrix expression in compiled array operations. The
single-window workload has different amortisation and allocation
costs. The thesis reports the delivered interface instead of
substituting a more favourable function with different semantics.

An additional professor-style aggregate experiment contains 84 paired
cases: seven source days, four sample fractions, and three repetitions.
Both backends converge in 79 pairs and both reach the iteration limit
in five. Summed likelihood-computation time is approximately 21.36~s
for LM and 24.33~s for SciPy; summed wall time is approximately
30.79~s and 33.76~s. This experiment uses its own initialisation and
a pooled stopping objective. It is evidence for that aggregate
workload and cannot replace production row-wise timings or be
described as an exact reproduction of the paper.

\begin{table}[tbp]
\centering\small
\begin{tabular}{@{}lrrrr@{}}
\toprule
Backend & Fits & Converged & Likelihood sum (s) & Wall-time sum (s) \\
\midrule
\TableRows{generated/aggregate}
\bottomrule
\end{tabular}
\caption{Aggregate numerical experiment. Each backend has 84 cases; the sums
include the recorded non-convergent cases. This is a different workload from
the production benchmark, not another estimate of the same median.}
\label{tab:aggregate}
\end{table}

Numerical stress tests and fallback address another dimension:
loss of significance when large gamma terms are subtracted. Their
purpose is to keep a finite complete score where ordinary
floating-point subtraction can be inaccurate. They do not add attack
information. In the measured detection workloads, backend agreement
provides no evidence that LM itself increases attack recall. Its
contribution is a working numerical integration with a measured
cost and an explicit correctness boundary.

\subsection{Measured effect of reusing native setup}

The historical backend change record supplies a direct comparison
of the earlier wrapper with the matrix entry point used here.
Its paired integration measurement reports a reduction from
114.68 to 55.58~$\mu$s for the 846-row likelihood workload and
from 0.30410 to 0.19486~s for complete fitting. These measurements
are recorded in \path{docs/lm_backend_change_record.tex}; they are
not a fresh experiment performed during thesis writing. The
separate production medians in Table~\ref{tab:backend-timing}
have their own measurement boundary and slightly different values.

The change record also reports bit-identical likelihood results
between the old and new native paths on 1,000 seeded stress
matrices. This supports the implementation claim that moving
shared setup outside the row loop need not change row-wise
semantics. It does not make the method bit-identical to SciPy,
whose numerical comparison is given separately.

The before/after improvement and the delivered SciPy comparison
answer different questions. Removing redundant native setup
substantially reduces the integration's overhead. The resulting
LM matrix kernel can nevertheless remain slower than SciPy on
the measured production workload. Reporting both comparisons
avoids attributing a wrapper improvement to a universal advantage
of the underlying mathematical method. It also shows that
backend selection depends on integration quality as well as
the approximation being evaluated.

\section{Window duration, common support, and response delay}
\label{sec:duration-results}

\subsection{Coarser windows can improve their own window metric}

The matched-duration study refits and recalibrates at 10, 20, 30, and
60 minutes. Table~\ref{tab:duration} gives static development results
on common observed support; Figure~\ref{fig:duration} separates F1
from delay. Samsung's ten-minute F1 here is 78.33\%, slightly
different from the reference 78.69\%, because the study uses its own
common-support fit/calibration preparation. The records should not
be merged as identical repetitions.

\begin{table}[tbp]
\centering\small
\begin{tabular}{@{}lrrrrrr@{}}
\toprule
Device & Min. & F1 & Intervals & By end & FP/24h & Delay (min) \\
\midrule
\TableRows{generated/duration}
\bottomrule
\end{tabular}
\caption{Static matched-duration development results. Intervals gives detected
over fully observed attack intervals; By end counts alerts available by attack
end. Median delay is conditional on detection. FP/24h uses observed exposure,
not benign-window count.}
\label{tab:duration}
\end{table}

\begin{figure}[tbp]
\centering
\begin{tikzpicture}
\begin{groupplot}[group style={group size=1 by 2,vertical sep=1.3cm},
 width=.92\textwidth,height=6cm,xmin=8,xmax=62,xtick={10,20,30,60},
 xlabel={Window duration (min)},ymajorgrids=true,grid style={gray!20},
 legend style={at={(.5,1.03)},anchor=south,legend columns=2,column sep=6pt}]
\nextgroupplot[ylabel={Window F1 (\%)},ymin=50,ymax=100]
\addplot[ChromecastBlue,mark=*,thick] table[x=minutes,y=f1,col sep=comma]{generated/duration_UNSW-Chromecast.csv};
\addplot[SamsungOrange,mark=square*,thick,dashed] table[x=minutes,y=f1,col sep=comma]{generated/duration_UNSW-Samsung-Camera.csv};
\legend{Chromecast,Samsung camera}
\nextgroupplot[ylabel={Median alert delay (min)},ymin=0,ymax=40]
\addplot[ChromecastBlue,mark=*,thick] table[x=minutes,y=delay,col sep=comma]{generated/duration_UNSW-Chromecast.csv};
\addplot[SamsungOrange,mark=square*,thick,dashed] table[x=minutes,y=delay,col sep=comma]{generated/duration_UNSW-Samsung-Camera.csv};
\end{groupplot}
\end{tikzpicture}
\caption{Static development results on common observed support, with fresh
fitting and calibration at every duration. Window F1 uses that duration's labels.
Delay runs from attack onset to first overlapping alert-window close and is
conditional on detection. Lines connect four measured settings.}
\label{fig:duration}
\end{figure}


Chromecast's window F1 rises from 61.22\% at ten minutes to 71.43\%
at one hour. Samsung reaches its highest value, 91.43\%, at thirty
minutes, then falls to 85.11\% at one hour. Pooling the corresponding
confusion counts gives about 73.37\%, 78.33\%, 85.11\%, and 81.97\%
for the four durations. This ranking describes duration-specific
tasks: the number and extent of positive windows change with scale.

The interval measure provides a different view. Chromecast detects
10, 9, 10, and 11 of 18 development attack intervals. Samsung detects
28, 28, 35, and 34 of 36. More detected intervals at a coarser scale
can help retrospective identification. It does not establish that an
alert arrives while the attack is active. Only one Chromecast
interval and five Samsung intervals receive an alert by attack end
at hourly resolution, compared with nine and 23 at ten minutes.

\subsection{A gain in coarse labels can conceal slower alerts}

Median alert delay rises from 6.83 to 36.93 minutes for Chromecast
and from 5.20 to 30.93 minutes for Samsung as duration increases
from ten minutes to one hour. Delay includes waiting for the first
alerting aggregate to close. Missed intervals are excluded from the
median, so these figures must be read together with interval recall.
A fast median over a small detected subset would not imply good
coverage.

The common ten-minute projection in
Table~\ref{tab:common-duration} avoids comparing only each duration's
own positive-window count. It assigns a coarser decision to its
covered ten-minute intervals on the same observed support. Under
this retrospective projection, ten minutes has the highest F1 for
both devices. Projection does not make decisions available earlier:
an hourly decision is still produced at the hourly close.

\begin{table}[tbp]
\centering\small
\begin{tabular}{@{}lrr@{}}
\toprule
Device & Original duration (min) & F1 on common ten-minute grid \\
\midrule
\TableRows{generated/duration_common_grid}
\bottomrule
\end{tabular}
\caption{Static development decisions projected onto a common ten-minute grid.
Projection permits a common labelling scale without altering alert availability.}
\label{tab:common-duration}
\end{table}

The later audit shows why a small positive denominator can make a
coarse result look particularly favourable. Hourly Chromecast F1
reaches 100\% on only two positive hourly windows, yet no alert is
available before the associated attacks end. This correctly describes
those hourly labels, but is insufficient evidence for timely response.
Duration is part of the detection task rather than a harmless
preprocessing setting.

\subsection{Adaptation remains vulnerable at other durations}

Changing duration changes the number of observations available to
calibrate tails and admit updates. The duration study scales nominal
buffers and cadence by time, but keeps the minimum calibration
sample requirement. Only 37 hourly calibration windows are available
for Samsung. A one-percent empirical tail at that sample size
depends on very few extreme observations.

Table~\ref{tab:duration-adaptation} gives the clearest adverse outcome.
At thirty minutes, Samsung threshold adaptation lowers development
F1 from 91.43\% to 83.33\% and increases false positives from two
to thirteen. On the later audit, false positives increase from zero
to 66 and F1 falls from 92.31\% to 17.50\%. At one hour, the
later-audit false positive count rises from three to 27.

\begin{table}[tbp]
\centering\small
\begin{tabular}{@{}rlrrrr@{}}
\toprule
Min. & Partition & Static F1 & Adaptive F1 & Static FP & Adaptive FP \\
\midrule
\TableRows{generated/duration_adaptation}
\bottomrule
\end{tabular}
\caption{Samsung threshold-adaptation failures in the matched-duration study.
Adaptive denotes its conservative threshold mode; the later partition is an
inspected audit, not a fresh selection holdout.}
\label{tab:duration-adaptation}
\end{table}

These results prevent a claim that the conservative label guarantees
non-degrading adaptation. The rule checks individual proposed changes;
it does not impose a global cumulative bound on every future threshold
trajectory. Small calibration sets and changes in accepted buffer
composition remain consequential. Periodic refitting promotes no
candidate in the duration study, so its static-equivalent outcomes
reflect rejection rather than an effective learned update.

Ten minutes is retained as the validated reference with the desired
response granularity. The evidence does not identify it as a universal
optimum. A deployment prioritising retrospective summaries could make
a different tradeoff, but would need to repeat calibration, delay
analysis, and adaptation checks at that duration.

\section{Closing final-partition evaluation}
\label{sec:final-results}

\subsection{Fixed settings and identical conservative modes}

The closing replay starts from the saved model/calibration pairs and
uses fixed ten-minute settings on 432 Chromecast and 576 Samsung
windows. Table~\ref{tab:final} gives its confusion counts. Static,
conservative threshold, and conservative refitting modes produce
identical binary predictions. Across the modes there are 3,024 saved
decisions, but only 1,008 unique labelled observations. The replays
must not be counted as three independent test samples.

\begin{table}[tbp]
\centering\small
\begin{tabular}{@{}lrrrrrrr@{}}
\toprule
Device & TP & FP & FN & TN & Precision & Recall & F1 \\
\midrule
\TableRows{generated/final}
\bottomrule
\end{tabular}
\caption{Closing final-partition results for each of the three identical
binary-decision modes. Pooled metrics use summed counts, not averaged device
metrics. The partition had prior retrospective label exposure; fixed replay
settings do not restore independence.}
\label{tab:final}
\end{table}

Chromecast has four true positives, three false positives, and two
false negatives, giving 57.14\% precision and 61.54\% F1. Samsung
has ten true positives, no false positives, and two false negatives,
giving 100\% precision and 90.91\% F1. Pooled precision is 82.35\%,
recall 77.78\%, and F1 80.00\%. The pooled false positive rate is
$3/990\simeq0.30\%$. Zero Samsung false positives is an observation
over 564 benign windows, not a bound on future false alerts.

The pooled F1 alone is insufficient. There are only six positive
Chromecast windows and twelve positive Samsung windows, while their
precision values differ sharply. Pooled 80.00\% does not mean that
a typical device has 80\% F1 or that the monitor will maintain it
under a different device mix. Individual counts preserve the
relevant denominators.

No threshold update or refit is promoted in the closing replay.
Threshold candidates fail acceptance checks, and refit attempts
are prevented by insufficient accepted data or rejected movement.
The initial pairs remain active. Equality of the modes is expected
and provides no evidence of an adaptation gain. It establishes
that the guarded modes can decline to change state.

The earlier replay records include strong rank discrimination on
these partitions, but absolute values depend on the saved score
version. Current count-first events can stop after a sufficient
count violation, so their score need not be the former full-feature
minimum. The closing conclusions use preserved decisions and interval
evidence rather than combining historical AUC values with the
current event definition.

\subsection{Attack intervals and alert availability}

Table~\ref{tab:final-intervals} expands the nine fully observed official
attack intervals. Eight are detected by at least one overlapping
positive window. Seven receive a first alert by attack end. The
distinction matters for the low-rate Samsung reflection interval:
it is counted as detected, but the first alert arrives 731~s after
onset, after the roughly ten-minute attack has ended.

\begin{table}[tbp]
\centering\footnotesize
\begin{tabular}{@{}l>{\raggedright\arraybackslash}p{5.2cm}rrr@{}}
\toprule
Device & Official attack type & Detected & Delay (s) & By end \\
\midrule
\TableRows{generated/final_intervals}
\bottomrule
\end{tabular}
\caption{Closing replay interval results. Attack names preserve annotation
vocabulary; they are not predicted attack types. Delay is measured to the first
overlapping alert-window close. A dash denotes no detection.}
\label{tab:final-intervals}
\end{table}

Chromecast misses \path{TcpSynReflection1W2D2W}, while detecting its
SSDP interval and \path{TcpSynReflection100W2D2W}. This is consistent
with weaker count evidence in the missed interval, but three cases
cannot establish a general relationship between the annotation's
rate class and detection probability. The detector sees counts; it
neither reads the rate label nor classifies the attack type.

All six Samsung intervals are detected, despite two positive windows
being missed. An interval can overlap several windows, so these
observations are compatible. Conditional median first-alert delay
is 414.5~s for Chromecast's two detected intervals and 448.5~s for
Samsung's six. These values are predominantly on the scale of window
closure, rather than the scoring kernel's millisecond scale.

The closing replay makes no new setting selection during execution.
Nevertheless, the same later labels were available in previous
scoring and duration audits. Its value is a fixed, fingerprinted
record of the selected implementation's behaviour on that partition.
It should not be presented as untouched external validation.

\subsection{Inspecting individual decisions}
\label{sec:event-cases}

Confusion counts establish how often predictions agree with annotations;
the saved event fields explain which statistical condition produced a
particular decision. Table~\ref{tab:event-cases} gives five selected
closing-replay examples. Cases A--C are Chromecast windows, and D--E
are Samsung windows. They are illustrative selections from the saved
events, not a new sample or a basis for recomputing aggregate rates.

\begin{table}[tbp]
\centering\footnotesize
\setlength{\tabcolsep}{4pt}
\begin{tabular}{@{}lrrrrrrrrl@{}}
\toprule
Case & TCP & UDP & SSDP & ARP & DM log $P$ & $m_D$ & $m_N$ & $m_A$ & Outcome \\
\midrule
\TableRows{generated/event_cases}
\bottomrule
\end{tabular}
\caption{Selected historical full-score events. A: Chromecast, 25 October
2018 at 07:20 UTC; B: 25 October at 07:50; C: 24 October at 16:30.
D: Samsung, 5 June 2018 at 18:50 UTC; E: 5 June at 19:00. Times are
window starts, and each window lasts ten minutes. Margins below zero
indicate a violated condition. Outcomes use the official window labels.}
\label{tab:event-cases}
\end{table}

Case A contains 267 SSDP packets alongside 188 TCP, 1,806 ordinary
UDP, and 41 ARP packets. Chromecast's fitted SSDP share is very small.
Its complete composition log probability is $-26.8698$, below the
saved composition cutoff of $-25.6108$. The composition margin is
$-1.311$, while the total and ARP margins are positive. This is a
composition-triggered alert rather than a total-volume violation.
The window overlaps the SSDP attack beginning at 07:25:23; its
07:30 close gives the recorded 277-second first-alert delay.

Case B is part of the missed low-rate reflection interval. It has
1,127 TCP packets and 1,118 ordinary UDP packets, no SSDP, and
51 ARP packets. Its complete score of $-25.5615$ is close to, but
above, the composition cutoff. The positive composition margin of
about $0.051$ and positive count margins make it non-alerting.
This is a false negative under the selected operating point, not
a numerical exception or a missing observation. The table does
not justify moving the cutoff to catch this one example: doing so
would require re-evaluating the benign tail and other attacks.

Case C demonstrates the opposite error. Its composition margin
is $-6.673$, giving an alert with no matching official attack
overlap. It is a false positive relative to the provided labels.
The four counts show the departure that the fitted rule sees,
but cannot establish whether it is a benign application change,
an unannotated event, or another cause. The event's UTC start
lies on 24 October even though it belongs to the configured
25 October capture ID. This is why evaluation uses actual UTC
bounds and capture context rather than assuming that nominal
capture dates are exact observation boundaries.

Case D violates both composition and total-volume conditions.
Its total is 2,068 packets, with a total margin of $-5.806$ and
a composition margin of $-19.353$. The historical full-feature
minimum is the latter. Under the current count-first order, the
negative total margin already establishes an alert, so composition
can be left unevaluated. The decision remains a true positive
even though the reported score and explanation fields can differ.
This concrete case illustrates the score-version boundary discussed
in Section~\ref{sec:joint-results}.

Case E follows that alerting Samsung window. Its counts are much
smaller, with total 517, and all margins are positive. The official
attack ends at 19:00:33, so the 19:00--19:10 window overlaps it for
33 seconds and is labelled positive. Its false-negative window
outcome is compatible with detection of the same interval by the
previous window. Neither the event schema nor the interval metric
is inconsistent; they answer different questions about the same
time range.

The examples show what interpretability means in this implementation.
A reader can recover the count vector, active reference, violated
margin, and annotation alignment. The explanation is local to the
statistical rule. It does not recover packet contents, prove
malicious intent, or assign an attack family. Preserving those
limits makes the event evidence useful without turning a numerical
score into a causal security diagnosis.

\section{Constrained Docker simulation}
\label{sec:edge-results}

\subsection{Processing cost includes rejected learning attempts}

The supplied-window runtime is measured using 2,000 and 20,000 input
records under one-CPU, 512~MiB container limits.
Table~\ref{tab:edge} reports both runs. Four of every five records
are observed windows and the fifth is missing-state input, giving
1,600 and 16,000 scored windows. The repeated fixture tests the
continuing process and update path; it is not a new labelled dataset.

\begin{table}[tbp]
\centering\small
\begin{tabular}{@{}lrr@{}}
\toprule
Measurement & Short replay & Long replay \\
\midrule
\TableRows{generated/edge}
\bottomrule
\end{tabular}
\caption{Supplied-window simulation measurements. Memory and disk units use
powers of two. Process high-water RSS and sparse sampled RSS are different
counters. No update becomes active in either run.}
\label{tab:edge}
\end{table}

Mean processing time is about 26.75~ms and 29.97~ms per scored
window. The worst windows take 4.449~s and 5.079~s. These tail
events are associated with attempted estimation; the longest
estimator calls are 4.430~s and 5.069~s. The long replay records
110 estimator calls and no promoted candidate. Rejecting an update
does not avoid paying its fitting cost.

The long history has an initial insufficient-data attempt followed
by 110 non-convergent refit attempts. The synthetic repeated counts
are not favourable evidence of successful model evolution. They do
exercise a repeated rejection path and show that the process remains
able to handle subsequent records without activating an invalid
candidate.

Scoring latencies are much smaller than a ten-minute observation
interval on this host. This supports the feasibility of processing
supplied completed windows. It cannot establish end-to-end response
time: collection, aggregation, transport, real-load queuing, and
waiting for window close are outside the processing timer.
Accelerated replay also keeps the allocated CPU busy; this is not
the expected utilisation when receiving one window every ten minutes.

\subsection{Finite memory and disk behaviour}

Process high-water RSS is 67.88~MiB in the short replay and 68.62~MiB
in the long replay. The long run's sparse samples reach 68.92~MiB,
with the last recorded samples stable at that level.
Figure~\ref{fig:memory} plots them. These measures remain separate
because they use different interfaces and sampling boundaries;
their discrepancy is not resolved by choosing the lower value.

\begin{figure}[tbp]
\centering
\begin{tikzpicture}
\begin{axis}[width=.95\textwidth,height=7cm,xlabel={Input windows processed},
 ylabel={Sampled process RSS (MiB)},ymin=60,ymax=75,xmin=0,xmax=20000,
 scaled x ticks=false,xtick={0,5000,10000,15000,20000},
 ymajorgrids=true,grid style={gray!20}]
\addplot[ChromecastBlue,thick,mark=*] table[x=windows,y=rss,col sep=comma]{generated/edge_memory.csv};
\end{axis}
\end{tikzpicture}
\caption{Recorded RSS samples during the 20,000-input synthetic Docker replay.
Lines connect sparse samples; they do not expose unobserved peaks or prove
indefinite memory stability. The workload uses the D-Link initial state and
promotes no update. The container memory limit is 512~MiB.}
\label{fig:memory}
\end{figure}


Measured memory remains far below the 512~MiB limit. This is
consistent with bounded buffers, synchronous output, and limited
retained history. A finite replay does not prove the absence of
every slow leak. It also does not establish how many devices can
share the allocation: the study has one process and one initial
device model.

Retained state occupies about 6.78~KiB and 31.25~KiB at the endpoints;
application logs occupy about 10.84~KiB and 109.55~KiB. Initial
growth reflects accumulating bounded history rather than promoted
models. These sizes exclude the image, captures, and downstream
storage of emitted events. The long run emits about 11.72~MB of
event data, an external storage responsibility if retained.

Application and Docker logging have rotation settings. The normal
simulation does not reach their limits; a separate reduced-limit
check exercises rotation. The 20,000-input run itself therefore
does not demonstrate eviction at production limits. The state
protocol has failure-injection tests, but this environment does
not reproduce physical power loss or storage-controller faults.

\section{Verification and evidence consistency}
\label{sec:verification-results}

The closing verification records the test suite at its own source
boundary. Later scoring-order work expands that suite. During
completion of this thesis, all 643 current tests pass, including
available container-dependent checks, using the locally available
edge image. This is current software verification, not a claim
that every experiment was rerun with newly generated models.

Checks cover input rejection, count semantics, numerical agreement,
short-circuit decisions, update acceptance, state recovery, streaming,
and command interfaces. They support implemented invariants. They
do not establish empirical detection generalisation, which requires
independent traffic observations rather than unit tests.

The document generator recomputes displayed confusion-matrix metrics,
verifies closing input fingerprints, and records the hashes of its
sources. Three source files have later versions covered by the
scoring-order regression; the remaining frozen dependencies retain
their recorded bytes. No raw capture is reprocessed and no historical
measurement is replaced during generation. Appendix~\ref{app:reproducibility}
identifies the commands and artifacts separating report regeneration
from experiment repetition.
